\documentclass[a4paper, 11pt]{article}
\usepackage{graphicx}
\usepackage{amsmath,array}
%\usepackage{tablefootnote}
\usepackage{epstopdf}
\author{Yuxi Chen}
\title{Algorithm}
\begin{document}
\maketitle
\newpage

\section{Unit conversion}

A physical quantity contains two parts: the value times the unit. For example, the weight of a table is $w = 50 kg = 50,000 g$, where 50 and 50,000 are the values, and kg and g are the units. If we define a new unit wg, which is defined as $1 wg = 2.5 kg$, then the table weight is $w = 20 wg$. In general, a physical quantity can be expressed as $u = \bar{u} u^*$, where $\bar{u}$ is the value/number and $u^*$ is the unit.  

\subsection{Mass normalization}

With the CGS unit, the particle motion equation is
\begin{eqnarray}
\frac{dv}{dt} = \frac{q}{m}(E + \frac{v}{c} \times B) 
\label{eq:velocity_update}
\end{eqnarray}
In FLEKS, all quantities are normalized. The normalization should not introduce new constants. So, from eq.\ref{eq:velocity_update}, the following expression is obtaied
\begin{equation}
\frac{q_{cgs}^* E_{cgs}^* t_{cgs}^*}{m_{cgs}^* v_{cgs}^*} = 1.
\label{eq:E1}
\end{equation}
From Gauss's law $\nabla\cdot E = 4\pi \rho$, we obtain 
\begin{equation}
E_{cgs}^* = \rho_{cgs} ^*x_{cgs}^*,
\label{eq:E2}
\end{equation}
where $\rho_{cgs}^* = q_{cgs}^*/x_{cgs}^{*3}$ is the charge density instead of mass density. From eq.\ref{eq:E1} and eq.\ref{eq:E2}, we obtain 
\begin{equation}
q_{cgs}^*=v_{cgs}^* \sqrt{m_{cgs}^* x_{cgs}^*}.
\label{eq:q1}
\end{equation}

In the code, $\bar{q}/(\bar{m}\bar{c}) = 1$ is required for a proton, so the normalization units $q_{cgs}^*$ and $m_{cgs}^*$ must satisfy 
\begin{equation}
q_{cgs}^*/(m_{cgs}^*v_{cgs}^*) = q_{cgs,p}/(m_{cgs,p}c_{cgs}),     
\end{equation}
where $q_{cgs,p}$ and $m_{cgs,p}$ are proton charge and mass in CGS unit, respectively. With eq.\ref{eq:q1}, we obtain the normalization mass in CGS unit:
\begin{eqnarray}
m_{cgs}^* = x_{cgs}^*(\frac{c_{cgs}m_{cgs,p}}{q_{cgs,p}})^2.
\end{eqnarray}
We note that $m_{cgs}^*$ also consists of the \textit{value} and the unit \textit{g}, and the expression above calculates its \textit{value}. Its \textit{value} in SI unit is calculated inside BATSRUS::ModPIC.f90, where the proton charge and mass are known in SI units. We need to convert the expression above into SI unit. Again, $m_{cgs}^*$ and $m_{SI}^*$ are the \textit{values} of the normalization mass in CGS and SI unit, respectively. If the normalization mass is 1.5kg, then $m_{SI}^* = 1.5$ and $m_{SI}^* = 1500$. The conversion is $m_{cgs}^* = 1000m_{SI}^*$. For the charge unit, $1C = 3\times 10^9 esu$ assume the speed of light is $3 \times 10^8m/s$. So, we have the following expression:
\begin{eqnarray}
1000m_{SI}^* = 100*x_{SI}^*(\frac{100*c_{SI}*1000*m_{SI,p}}{3\times10^9 q_{SI,p}})^2 \\
m_{SI}^* = 10^7x_{SI}^*(\frac{m_{SI,p}}{q_{SI,p}})^2.
\end{eqnarray}

Once $m_{SI}^*$ is obtained, it is passed to FLEKS and converted to $m_{cgs}^*$. 

\subsection{Length and velocity normalization}

These two are free parameters. 

\subsection{Charge and current normalization}

$q_{cgs}^*=v_{cgs}^* \sqrt{m_{cgs}^* x_{cgs}^*}$, as it is shown in eq.\ref{eq:q1}. $j_{cgs}^* = q_{cgs}^*v_{cgs}^*/x_{cgs}^{*3}$

\subsection{B and E normalization}

B and E have the same unit in CGS. Substituting eq.\ref{eq:q1} into eq.\ref{eq:E1} to obtain $E_{cgs}^* = B_{cgs}^* = \sqrt{\frac{m_{cgs}^*}{x_{cgs}^*3}}v_{cgs}^*$ 

\subsection{The consequences of changing light speed}

Reducing the light speed changes the EM wave speed in the simulation, but it does change the interaction between magnetic field and particles. When we discuss the mass normalization, $\bar{q}/(\bar{m}\bar{c}) = 1$ instead of $\bar{q}/\bar{m} = 1$ is required, because there is no assumption of $\bar{c}=1$. The light speed in eq.\ref{eq:velocity_update} is always the physical light speed no matter what $v_{cgs}^*$ is choosen. For example, if $v_{cgs}^* = 0.1c$, $\bar{c}$ has to be 10, otherwise, $\bar{q}/(\bar{m}\bar{c}) = 1$ can not be satisfied. 

Since reducing $v_{cgs}^*$ does not change the interaction between magnetic field and particle, the particle gyromotion, including gyro-radius and gyro-frequency, does not change. The inertial length is the gyro-radius of a particle with Alfven velocity, and neither gyromotion nor Alfven velocity changes, so the inertial length does not change with $v_{cgs}^*$.

How about the interaction between the electric field and particles? Since $\bar{c} = \bar{q}/\bar{m}$, which is not necessarily be 1, the Coulomb force for a proton should be $\bar{q}/\bar{m}E = \bar{c}E$. However, in the code, $\bar{c}$ is ignored, which suggests the Coulomb force in the simulation is $\bar{c}$ times weaker than in reality. Is this reasonable? What is its consequence? It is to be clarified. 

Most PIC simulations use reduced light speed, maybe the results are interpreted in a different way than FLEKS/MHD-EPIC, but I believe they also change the ratio between the Coulomb force and the $v \times B$ force. So, our approach of reducing light speed is probably reasonable. 

\section{Particle mover}

\subsection{Boris algorithm}

FLEKS uses Boris particle mover: 
\begin{eqnarray}
\frac{\mathbf{v}^{n+1} - \mathbf{v}^{n}}{\Delta t} = \frac{q}{m}(\mathbf{E}^{n+\theta} + \bar{\mathbf{v}}\times\mathbf{B}) \\
\frac{\mathbf{x}^{n+1/2} - \mathbf{x}^{n-1/2}}{\Delta t} = \mathbf{v}^n
\end{eqnarray}
where $\mathbf{v}$ is defined as $\bar{\mathbf{v}} = \frac{\mathbf{v}^{n+1} + \mathbf{v}^{n}}{2}$. The traditional method updates the velocity with the following steps.

\begin{enumerate}
    \item Acceleration 
\begin{eqnarray}
    \mathbf{v}^- = \mathbf{v}^{n} + \frac{q\Delta t}{2m}\mathbf{E}^{n+\theta}
\end{eqnarray}

    \item Rotation
\begin{eqnarray} 
    \mathbf{a} = \mathbf{v}^- + \mathbf{v}^- \times \mathbf{t} \\
    \mathbf{v}^+ = \mathbf{v}^- + \mathbf{a} \times \mathbf{s}
\end{eqnarray}
    where $\mathbf{t} = \frac{q\Delta t}{2m} \mathbf{B}$, and $\mathbf{s} = \frac{2\mathbf{t}}{1+t^2}$.

    \item  Acceleration
\begin{eqnarray}    
    \mathbf{v}^{n+1} = \mathbf{v}^+ + \frac{q\Delta t}{2m}\mathbf{E}^{n+\theta}
\end{eqnarray}
\end{enumerate}

FLEKS solves the same velocity equation, but the implementation is slightly different.
\begin{enumerate}
    \item Acceleration 
\begin{eqnarray}
    \mathbf{v}^- = \mathbf{v}^{n} + \frac{q\Delta t}{2m}\mathbf{E}^{n+\theta}
\end{eqnarray}

\item Calculate $\bar{\mathbf{v}}$. 

With the expressions of $\mathbf{v}^+$ and $\mathbf{v}^-$ above, it is easy (it is probably not straightforwad, but I do not know how to number equations in markdown) to obtain $\mathbf{v}^{n+1} - \mathbf{v}^{n} = \mathbf{v}^+ - \mathbf{v}^- + \frac{q\Delta t}{m} \mathbf{E}^{n+\theta}$. Substitute it into the velocity update equation to remove electric field from the expression $\frac{\mathbf{v}^{+} - \mathbf{v}^{-}}{\Delta t} = \frac{q}{m}\bar{\mathbf{v}}\times\mathbf{B}$. Since $\mathbf{v}^+ = 2\bar{\mathbf{v}} - \mathbf{v}^-$, we obtain the equation for $\bar{\mathbf{v}}$:
\begin{eqnarray}
  \bar{\mathbf{v}} = \mathbf{v}^- + \frac{q\Delta t}{2m} \bar{\mathbf{v}} \times \mathbf{B} = \mathbf{v}^- + \bar{\mathbf{v}} \times \mathbf{t}   
\end{eqnarray}

To solve this equation, we apply $\cdot \mathbf{t}$ to the expression of $\bar{\mathbf{v}}$ to obtain $\bar{\mathbf{v}} \cdot \mathbf{t} = \mathbf{v}^- \cdot \mathbf{t}$. Then we apply $\times \mathbf{t}$ to the expression of $\bar{\mathbf{v}}$ to obtain
\begin{eqnarray}
\bar{\mathbf{v}} \times \mathbf{t} = \mathbf{v}^- \times \mathbf{t} + \bar{\mathbf{v}} \times \mathbf{t} \times \mathbf{t} = \mathbf{v}^- \times \mathbf{t} + (\bar{\mathbf{v}} \cdot \mathbf{t})\mathbf{t} - t^2\bar{\mathbf{v}} = \mathbf{v}^- \times \mathbf{t} + (\mathbf{v}^- \cdot \mathbf{t})\mathbf{t} - t^2\bar{\mathbf{v}}
\end{eqnarray}
Substitute the expression of $\bar{\mathbf{v}} \times \mathbf{t}$ to the expression of $\bar{\mathbf{v}}$ to obtain 
\begin{eqnarray}
\bar{\mathbf{v}} = \frac{\mathbf{v}^- + \mathbf{v}^- \times \mathbf{t} + (\mathbf{v}^- \cdot \mathbf{t})\mathbf{t}}{1+t^2}
\end{eqnarray}

\item Then it is easy to obtain $\mathbf{v}^{n+1} = 2*\bar{\mathbf{v}} - \mathbf{v}^{n}$
\end{enumerate}

\subsection{Relativistic Boris algorithm}

With $\mathbf{u} = \gamma \mathbf{v}$, the velocity update equation is
\begin{eqnarray}
\frac{\mathbf{u}^{n+1} - \mathbf{u}^{n}}{\Delta t} = \frac{q}{m}(\mathbf{E}^{n+\theta} + \bar{\mathbf{u}}\times \frac{\mathbf{B}}{\gamma ^{n+1/2}})
\end{eqnarray}

\begin{enumerate}
    \item Convert $\mathbf{v}$ to $\mathbf{u}$
    \item Acceleration
\begin{eqnarray}
    \mathbf{u}^- = \mathbf{u}^{n} + \frac{q\Delta t}{2m}\mathbf{E}^{n+\theta}
\end{eqnarray}

\item Calculate $\bar{\mathbf{u}}$ with
\begin{eqnarray}
\bar{\mathbf{u}} = \frac{\mathbf{u}^- + \mathbf{u}^- \times \mathbf{t} + (\mathbf{u}^- \cdot \mathbf{t})\mathbf{t}}{1+t^2}
\end{eqnarray}
where $\mathbf{t} = \frac{q\Delta t}{2m} \mathbf{B}/\gamma^{n+1/2}$. From the definition of $\gamma$ and $\mathbf{u} = \gamma \mathbf{v}$, $\gamma ^{n+1/2} = \sqrt{1+(u^-/c)^2}$ is obtained. 

\item Then it is easy to obtain $\mathbf{u}^{n+1} = 2*\bar{\mathbf{u}} - \mathbf{u}^{n}$
\item  Convert $\mathbf{u}$ back to $\mathbf{v}$
\end{enumerate}

\subsection{Boris mover linearization}

The Boris velocity update can be decomposed into an electric acceleration, a magnetic rotation, and a second electric acceleration. Denoting $\mathbf{t} = \frac{q\Delta t}{2m}\mathbf{B}^n$, the average velocity $\bar{\mathbf{v}} = \frac{\mathbf{v}^{n+1} + \mathbf{v}^n}{2}$ is related to the effective electric field $\mathbf{E}^{n+\theta}$ by:
\begin{equation}
\bar{\mathbf{v}} = \hat{\mathbf{v}} + \boldsymbol{\alpha}_p \cdot \mathbf{E}^{n+\theta},
\label{eq:vbar_linearized}
\end{equation}
where $\hat{\mathbf{v}}$ is the tentative velocity depending exclusively on the known velocity $\mathbf{v}^n$ and magnetic field $\mathbf{B}^n$:
\begin{equation}
\hat{\mathbf{v}} = \frac{\mathbf{v}^n + \mathbf{v}^n \times \mathbf{t} + (\mathbf{v}^n \cdot \mathbf{t})\mathbf{t}}{1 + t^2},
\end{equation}
and $\boldsymbol{\alpha}_p$ is the single-particle susceptibility tensor:
\begin{equation}
\boldsymbol{\alpha}_p = \frac{q\Delta t}{2m(1 + t^2)} \left[ \mathbf{I} + \mathbf{t}\mathbf{t} + \mathbf{I} \times \mathbf{t} \right].
\end{equation}
Here, $(\mathbf{I} \times \mathbf{t}) \cdot \mathbf{E} = \mathbf{E} \times \mathbf{t}$. Equation~\eqref{eq:vbar_linearized} provides the exact linear closure that couples particle kinetics to the implicit electromagnetic field equations in the full PIC semi-implicit method.

\section{Full PIC semi-implicit solver}

In the full-PIC approach, both electrons and ions are treated as kinetic particle species. To bypass the explicit Courant--Friedrichs--Lewy (CFL) constraint associated with the speed of light and plasma frequency ($\omega_{pe}\Delta t \le 2$), FLEKS employs the semi-implicit particle-in-cell method.

\subsection{Maxwell--Vlasov semi-implicit discretization}

Maxwell's equations in normalized units ($c=1$) are discretized in time using a $\theta$-implicit parameter ($\theta \in [0.5, 1.0]$):
\begin{align}
\frac{\mathbf{B}^{n+1} - \mathbf{B}^n}{\Delta t} &= -\nabla \times \mathbf{E}^{n+\theta}, \label{eq:maxwell_faraday} \\
\frac{\mathbf{E}^{n+1} - \mathbf{E}^n}{\Delta t} &= \nabla \times \mathbf{B}^{n+\theta} - 4\pi \bar{\mathbf{J}}, \label{eq:maxwell_ampere}
\end{align}
where intermediate fields are evaluated as:
\begin{align}
\mathbf{E}^{n+\theta} &= \theta \mathbf{E}^{n+1} + (1-\theta)\mathbf{E}^n, \\
\mathbf{B}^{n+\theta} &= \theta \mathbf{B}^{n+1} + (1-\theta)\mathbf{B}^n = \mathbf{B}^n - \theta \Delta t \nabla \times \mathbf{E}^{n+\theta}.
\end{align}
Substituting $\mathbf{B}^{n+\theta}$ into Amp\`ere's law~\eqref{eq:maxwell_ampere} yields:
\begin{equation}
\mathbf{E}^{n+\theta} - \mathbf{E}^n = \theta \Delta t \left[ \nabla \times \left(\mathbf{B}^n - \theta \Delta t \nabla \times \mathbf{E}^{n+\theta}\right) - 4\pi \bar{\mathbf{J}} \right].
\label{eq:ampere_intermediate}
\end{equation}

\subsection{Susceptibility mass matrix and tentative current}

Using the linearized Boris mover relation~\eqref{eq:vbar_linearized}, the average current density $\bar{\mathbf{J}}(\mathbf{x})$ is expressed as:
\begin{equation}
\bar{\mathbf{J}}(\mathbf{x}) = \hat{\mathbf{J}}(\mathbf{x}) + \mathbf{M} \cdot \mathbf{E}^{n+\theta},
\end{equation}
where $\hat{\mathbf{J}}(\mathbf{x})$ is the tentative current density deposited onto grid nodes using shape functions $S(\mathbf{x} - \mathbf{x}_p)$:
\begin{equation}
\hat{\mathbf{J}}(\mathbf{x}) = \sum_p q_p \hat{\mathbf{v}}_p S(\mathbf{x} - \mathbf{x}_p),
\end{equation}
and $\mathbf{M}$ is the nodal susceptibility mass matrix:
\begin{equation}
\mathbf{M}(\mathbf{x}, \mathbf{x}') = \sum_p q_p S(\mathbf{x} - \mathbf{x}_p) \boldsymbol{\alpha}_p S(\mathbf{x}' - \mathbf{x}_p).
\end{equation}
The mass matrix $\mathbf{M}$ is accumulated at node locations with support across neighboring stencil points, supporting block-adaptive mesh refinement (AMR).

\subsection{Implicit electric field linear system}

Substituting $\bar{\mathbf{J}}$ into eq.~\eqref{eq:ampere_intermediate} and rearranging gives the implicit linear system for $\mathbf{E}^{n+\theta}$:
\begin{equation}
\left[ \mathbf{I} + \theta^2 \Delta t^2 \nabla \times (\nabla \times \cdot) + 4\pi \theta \Delta t \mathbf{M} \right] \mathbf{E}^{n+\theta} = \mathbf{R},
\label{eq:esolve_operator}
\end{equation}
where the right-hand side $\mathbf{R}$ is evaluated from the known state at $t^n$:
\begin{equation}
\mathbf{R} = \mathbf{E}^n + \theta \Delta t \left( \nabla \times \mathbf{B}^n - 4\pi \hat{\mathbf{J}} \right).
\end{equation}
Using the vector identity $\nabla \times (\nabla \times \mathbf{E}) = \nabla (\nabla \cdot \mathbf{E}) - \nabla^2 \mathbf{E}$, the linear operator is implemented matrix-free and solved iteratively using GMRES or Jacobian-Free Newton--Krylov.

Once $\mathbf{E}^{n+\theta}$ is obtained, the fields are advanced:
\begin{align}
\mathbf{E}^{n+1} &= \frac{1}{\theta}\mathbf{E}^{n+\theta} - \frac{1-\theta}{\theta}\mathbf{E}^n, \\
\mathbf{B}^{n+1} &= \mathbf{B}^n - \Delta t \nabla \times \mathbf{E}^{n+\theta}.
\end{align}
To maintain Gauss's law $\nabla \cdot \mathbf{E} = 4\pi \rho_c$, an optional div($\mathbf{E}$) cleaning step solves the Poisson equation $\nabla^2 \phi = \frac{1}{\theta}(\nabla \cdot \mathbf{E} - 4\pi \rho_c)$ and applies a localized particle position correction $\Delta \mathbf{x} \propto -\nabla \phi$ or electric field projection.

\section{Upwind schemes and comoving frame solving}

In simulations featuring high-speed plasma flows, shocks, and sharp boundary layers, standard central-difference discretizations can produce unphysical dispersion, numerical oscillations, and grid heating. FLEKS incorporates a comoving frame formulation and directional upwind artificial viscosity schemes for both the electric and magnetic fields.

\subsection{Comoving frame field solving}

When plasmas possess large bulk drift velocities (such as in the supersonic solar wind), solving fields in the stationary laboratory frame Doppler-shifts high-frequency wave modes and requires resolving fast advection. FLEKS provides the option to solve the electromagnetic fields in the frame comoving with the plasma bulk flow:

\begin{enumerate}
  \item The background fluid velocity $\mathbf{U}_0$ is computed from the ion fluid moments deposited at grid nodes:
  \begin{equation}
  \mathbf{U}_0(\mathbf{x}) = \frac{\sum_s (\rho \mathbf{u})_s}{\sum_s \rho_s}.
  \end{equation}
  Optionally, $\mathbf{U}_0$ may be smoothed with a digital filter to suppress PIC shot noise.
  \item The motional electric field in the comoving frame is:
  \begin{equation}
  \mathbf{E}_0 = -\mathbf{U}_0 \times \mathbf{B}^n.
  \end{equation}
  \item The tentative particle velocities and tentative current $\hat{\mathbf{J}}$ are evaluated relative to the comoving frame: $\hat{\mathbf{v}}' = \hat{\mathbf{v}} - \mathbf{U}_0$.
  \item The motional electric field generates an effective source current through the susceptibility matrix. Its contribution $\mathbf{M} \cdot \mathbf{E}_0$ is added to the right-hand side of the implicit system~\eqref{eq:esolve_operator}:
  \begin{equation}
  \mathbf{R}_{\text{comov}} = \mathbf{R} + 4\pi \theta \Delta t \, \mathbf{M} \cdot \mathbf{E}_0.
  \end{equation}
\end{enumerate}

\subsection{Upwind scheme for electric field}

A directional Lax--Friedrichs artificial viscosity is added to the semi-implicit electric field equation to suppress numerical oscillations near shocks and steep gradients.

The artificial diffusion flux at cell interface $i+1/2$ along spatial direction $d$ is:
\begin{equation}
F_{i+1/2}^{\text{art}} = \frac{1}{2} c_R u_R (E_{i+1} - E_i),
\end{equation}
where $u_R$ is the characteristic propagation velocity. Two choices are supported:
\begin{itemize}
  \item $u_R = c_{\max}$, the (normalized) speed of light, giving isotropic maximum-speed upwinding.
  \item $u_R = |(\mathbf{U}_0)_d|$, the local component of the background plasma velocity along the interface direction. This limits numerical dissipation to the actual advection velocity and is used in conjunction with the comoving frame solver.
\end{itemize}
The coefficient $c_R$ is governed by a generalized Minmod/Sweby flux limiter:
\begin{equation}
c_R = 1 - \phi(r_R), \qquad \phi(r) = \max\left(0,\, \min\left(\theta_E r,\; \frac{1+r}{2},\; \theta_E\right)\right), \quad r = \frac{E_i - E_{i-1}}{E_{i+1} - E_i},
\end{equation}
where $\theta_E \in [0, 2]$ controls the amount of limiting:
\begin{itemize}
  \item $\theta_E = 0$: $\phi(r) = 0$, $c_R = 1$ --- fully diffusive first-order Lax--Friedrichs upwind scheme.
  \item $\theta_E = 1$: reproduces the minmod limiter.
  \item $\theta_E = 2$: least diffusive TVD limiter.
\end{itemize}

In the implicit electric field operator~\eqref{eq:esolve_operator}, this upwind term enters the left-hand side:
\begin{equation}
\Delta \mathcal{A}(\mathbf{E})_i = -\frac{\theta \Delta t}{2\Delta x} \left[ c_R u_R (E_{i+1} - E_i) - c_L u_L (E_i - E_{i-1}) \right].
\end{equation}

When the limiter coefficients $c_{R/L}$ are evaluated from the lagged field $\mathbf{E}^n$, the operator remains strictly linear and GMRES is used. When $c_{R/L}$ depends on the current iterate $\mathbf{E}$, the operator is non-linear and a matrix-free Newton--Krylov solver is applied instead.

\subsection{Upwind scheme for magnetic field}

A directional upwind diffusion correction is applied to the cell-centered magnetic field following the Faraday update. The correction applied to cell $i$ is:
\begin{equation}
\Delta B_i = \frac{\Delta t}{2\Delta x} \left[ c_R u_R (B_{i+1} - B_i) - c_L u_L (B_i - B_{i-1}) \right].
\end{equation}
Key properties of the upwind $\mathbf{B}$ scheme include:
\begin{itemize}
  \item \textbf{Advection velocity:} $u_R$ and $u_L$ are face-reconstructed from the background plasma velocity $\mathbf{U}_0$. A constant advection speed can alternatively be prescribed.
  \item \textbf{Slope limiter:} $c_{R/L} = 1 - \phi(r)$ uses the same Minmod/Sweby limiter family as the electric field scheme, with an independent parameter $\theta_B \in [0, 2]$.
  \item \textbf{Boundary protection:} At physical domain boundaries, diffusion is suppressed when the flow is directed outward, preventing spurious boundary reflections.
\end{itemize}

\section{Hybrid PIC solver}

In the hybrid PIC approach, ions are treated as kinetic particles and electrons are modeled as a massless, charge-neutralizing fluid. The electromagnetic fields are evolved through a generalized Ohm's law rather than through the full implicit Maxwell solver.

\subsection{The generalized Ohm's law}

The electric field is assembled cell-centered from the ion fluid moments and the magnetic field:
\begin{equation}
\mathbf{E} = -\mathbf{U}_i \times \mathbf{B}
           + \eta\,\mathbf{J}
           + \frac{\mathbf{J} \times \mathbf{B}}{\rho_q}
           - \frac{\nabla P_e}{\rho_q}
           - \frac{\eta_h}{4\pi}\,\nabla \times (\nabla^2 \mathbf{B}),
\label{eq:ohmslaw}
\end{equation}
where:
\begin{itemize}
  \item $-\mathbf{U}_i \times \mathbf{B}$ is the \emph{convection} term, always active.
  \item $\eta\,\mathbf{J}$ is the \emph{resistive} term, with $\eta$ the magnetic diffusivity.
  \item $(\mathbf{J}\times\mathbf{B})/\rho_q$ is the \emph{Hall} term. When the Hall term is absent, the Hall-whistler time-step restriction is removed.
  \item $-\nabla P_e/\rho_q$ is the \emph{electron pressure gradient} term. The electron pressure $P_e$ supports two closures:
  \begin{itemize}
    \item Isothermal ($\gamma=1$): $P_e = T_e \rho$.
    \item Adiabatic ($\gamma > 1$): $P_e = P_0 (\rho/\rho_0)^\gamma$, where $P_0 = \rho_0 T_e$.
  \end{itemize}
  A minimum density floor $\rho_{\min}$ is applied to the $1/\rho_q$ factors in the Hall and electron pressure terms to prevent unbounded contributions in near-vacuum cells.
  \item $-(\eta_h/4\pi)\,\nabla\times(\nabla^2\mathbf{B})$ is the \emph{hyper-resistivity} term. Hyper-resistivity damps grid-scale PIC shot noise without introducing excessive dissipation at large scales. The coefficient $\eta_h$ can be set to scale with the grid resolution:
  \begin{equation}
  \eta_h = 4\pi C_h \frac{\Delta x_{\min}^4}{\Delta t_{\text{sub}}},
  \end{equation}
  where $C_h$ is a dimensionless tuning parameter, typically $10^{-3}$--$10^{-2}$.
\end{itemize}
The current density is obtained from Amp\`ere's law (CGS), $\mathbf{J} = \nabla\times\mathbf{B}/(4\pi)$. The explicit resistive and hyper-resistive diffusion terms must satisfy parabolic CFL stability bounds:
\begin{equation}
\frac{\eta k_{\max}^2 \Delta t_{\text{sub}}}{4\pi} \le C_{\text{CFL}}, \qquad \frac{\eta_h k_{\max}^4 \Delta t_{\text{sub}}}{4\pi} \le C_{\text{CFL}},
\end{equation}
where $C_{\text{CFL}} = 2.785$ for RK4 and $2.513$ for SSPRK3.

\subsection{Cell-centered field layout}

All hybrid fields ($\mathbf{B}$, $\mathbf{E}$, $\mathbf{J}$, and ion moments $\rho$, $\rho\mathbf{u}$) are stored cell-centered. Spatial curl and gradient operators use a $2\,\Delta x$ central-difference stencil:
\begin{equation}
(\nabla \times \mathbf{B})_{i,j,k} \propto \frac{\mathbf{B}_{i+1} - \mathbf{B}_{i-1}}{2\Delta x},
\end{equation}
whose discrete Fourier kernel vanishes identically at the grid Nyquist mode ($k\Delta x = \pi$). This eliminates the grid-scale odd-even decoupling instability that would otherwise force heavy Hall sub-cycling, and removes the need for center$\leftrightarrow$node interpolations during field advances.

\subsection{Time integration and subcycling of $\mathbf{B}$}

The magnetic field is advanced via Faraday's law:
\begin{equation}
\frac{\partial \mathbf{B}}{\partial t} = -\nabla\times\mathbf{E},
\end{equation}
using an explicit Runge--Kutta integrator. Two choices are provided:
\begin{itemize}
  \item \textbf{RK4}: Classical fourth-order Runge--Kutta. Because RK4 has a wide stability region along the imaginary axis (where Hall-whistler oscillations reside), a single subcycle per coarse step is often sufficient.
  \item \textbf{SSPRK3}: Strong-stability-preserving third-order Runge--Kutta with time-centered intermediate stages.
\end{itemize}
Within one coarse particle time step $\Delta t$, $N_{\text{sub}}$ subcycles of step size $\Delta t_{\text{sub}} = \Delta t / N_{\text{sub}}$ are taken.

To achieve second-order temporal accuracy, the ion moments are linearly interpolated across each magnetic sub-step fraction $h_{\text{step}} \in [0, 1]$ between the previous deposit $X^{n-1/2}$ and the current deposit $X^{n+1/2}$:
\begin{equation}
X(h_{\text{step}}) = \left(\tfrac{1}{2}-h_{\text{step}}\right) X^{n-1/2} + \left(\tfrac{1}{2}+h_{\text{step}}\right) X^{n+1/2}.
\end{equation}

\subsection{Algorithmic cycle and noise mitigation}

The complete hybrid PIC time step proceeds through the following stages:
\begin{enumerate}
  \item \textbf{Particle Boris push:} Ions are advanced $\mathbf{v}^n \to \mathbf{v}^{n+1}$ using integer-step fields $\mathbf{E}^n, \mathbf{B}^n$; positions are updated as $\mathbf{x}^{n+1/2} = \mathbf{x}^{n-1/2} + \Delta t\,\mathbf{v}^n$.
  \item \textbf{Kinetic source and boundary updates:} Particle boundary conditions, chemistry, charge exchange, and physical inflow flux injection are evaluated.
  \item \textbf{Moment deposition:} Ion charge $\rho_q^{n+1/2}$ and mass momentum $(\rho\mathbf{u})_i^{n+1/2}$ are deposited cell-centered. A digital-filter smoothing can be applied to reduce particle noise.
  \item \textbf{Faraday advance:} Subcycled explicit RK4 or SSPRK3 advances $\mathbf{B}^n \to \mathbf{B}^{n+1}$ using the interpolated ion moments.
  \item \textbf{Electric field assembly:} Cell-centered $\mathbf{E}^{n+1}$ is computed from eq.~\eqref{eq:ohmslaw} and stored for the subsequent Boris push.
\end{enumerate}

To suppress high-frequency PIC shot noise in the Hall term and particle trajectories, an optional exponential moving average of the magnetic field is maintained:
\begin{equation}
\mathbf{B}_{\text{avg}}^{n+1} = \alpha \mathbf{B}_{\text{avg}}^n + (1-\alpha)\mathbf{B}^{n+1}, \qquad \alpha = 1 - \frac{1}{N_{\text{avg}}}.
\end{equation}
When active, $\mathbf{B}_{\text{avg}}$ replaces the instantaneous $\mathbf{B}$ inside the convection and Hall terms of Ohm's law and in the Boris particle pusher.

\section{Miscellaneous}
\subsection{Calculate total pressure tensor from sub groups}
Pressure tensor for each sub group is:
\begin{eqnarray}
    p^s_{m,n} = \frac{1}{V} \sum_{i=1} w_i^s (v_{i,m}^s-\bar{v}_m^s) (v_{i,n}^s-\bar{v}_n^s) 
\end{eqnarray}
where $V$ is the volume, and $\bar{v}_n^2$ and $\bar{v}_m^s$ are the average velocity of the sub group $s$. 

The total pressure tensor is:
\begin{eqnarray}
    p_{m,n} = \frac{1}{V} \sum_{s}\sum_{i=1} w_i^s (v_{i,m}^s-\bar{v}_m) (v_{i,n}^s-\bar{v}_n)\\
    \bar{v}_m = \frac{\sum_{s}\sum_{i=1} w_i^s v_{i,m}^s}{\sum_{s}\sum_{i=1} w_i^s} = \frac{\sum_s \rho_s \bar{v}_m}{\sum_{s}\rho_s} \\
    \bar{v}_n = \frac{\sum_{s}\sum_{i=1} w_i^s v_{i,n}^s}{\sum_{s}\sum_{i=1} w_i^s} = = \frac{\sum_s \rho_s \bar{v}_n}{\sum_{s}\rho_s}
\end{eqnarray}
where $\rho_s$ is the density of sub group $s$, $\bar{v}_n$ and $\bar{v}_m$ are the average velocity of all particles

\begin{eqnarray}
    p_{m,n} = \frac{1}{V} \sum_{s}\sum_{i=1} w_i^s [(v_{i,m}^s- \bar{v}_m^s) + (\bar{v}_m^s -\bar{v}_m)] [(v_{i,n}^s- \bar{v}_n^s) + (\bar{v}_n^s  - \bar{v}_n)]\\
    = \frac{1}{V} \sum_{s}\sum_{i=1} w_i^s[(v_{i,m}^s- \bar{v}_m^s)(v_{i,n}^s- \bar{v}_n^s)] + \\ \frac{1}{V} \sum_{s}\sum_{i=1} w_i^s[(v_{i,m}^s- \bar{v}_m^s)(\bar{v}_n^s - \bar{v}_n)] + \\
    \frac{1}{V} \sum_{s}\sum_{i=1} w_i^s[(\bar{v}_m^s -\bar{v}_m)(v_{i,n}^s- \bar{v}_n^s)] + \\ \frac{1}{V} \sum_{s}\sum_{i=1} w_i^s[(\bar{v}_m^s -\bar{v}_m)(\bar{v}_n^s - \bar{v}_n)]
\end{eqnarray}
The second and third terms are zeros. The first term is simply the sum of all the sub group tensors. The fourth term is $\sum_s \rho_s [(\bar{v}_m^s -\bar{v}_m)(\bar{v}_n^s - \bar{v}_n)]$. So, the final expression is:
\begin{eqnarray}
    p_{m,n} = \sum_s p^s_{m,n} + \sum_s \rho_s [(\bar{v}_m^s -\bar{v}_m)(\bar{v}_n^s - \bar{v}_n)]
\end{eqnarray}
\end{document}